\section{USAMO 1999}
        \begin{enumerate}
        	\item (\textit{USAMO 1999}) Some checkers placed on an $n\times n$ checkerboard satisfy the following conditions:


 (a) every square that does not contain a checker shares a side with one that does;


 (b) given any pair of squares that contain checkers, there is a sequence of squares containing checkers, starting and ending with the given squares, such that every two consecutive squares of the sequence share a side.


 Prove that at least $(n^{2}-2)/3$ checkers have been placed on the board.


 



\textbf{Solution 1}

For the proof let's look at the checkers board as a graph $G$, where the checkers are the vertices and the edges are every pair of checkers sharing a side.

 Define $R$ as the circuit rank of $G$(the minimum number of edges to remove from a graph to remove all its cycles).

 Define $x$ as the number of checkers placed on the board.

 1. From (a) for every square containing a checker there can be up to 4 distinct squares without checkers; thus, $4\times x + x = 5\times x \ge n^2$ (this is the most naive upper bound).

 2. From (b) this graph has to be connected. which means that no square which contains a checker can share 4 sides with squares that doesn't contain a checker (because it has to share at least one side with a square that  contains a checker).

 3. Now it is possible to improve the upper bound. Since every edge in G represents a square with a checker that shares a side with a square with another checker, the new upper bound is $5\times x - \text{[sum of degrees in G]} \ge n^2.$ 

 4. Thus, $5x - (2x - 2 + 2R) \ge n^2$ (see lemma below).

 5. Thus,  $3x + 2 - 2R \ge n^2.$

 6. Thus,  $x \ge \frac{n^2 - 2 + 2R}3,$ where $R\ge0$ and is equal to 0 if $G$ is a tree.

 $\textit{Clarification/Lemma:}$ The sum of degrees of a connected graph $G = (V,E)$ is $2V -2 + 2R = 2E,$ where $R$ is the circuit rank of $G$.

 $\textit{Proof.}$ $2 \times V -2$ is the sum of ranks of the spanning tree created by decircuiting the graph $G$. Since the circuit rank of $G$ is $R$, $2R$ is the sum of ranks removed from $G$. Thus,  $2V -2 + 2R = 2 E.$ $\blacksquare$

 



\textbf{Solution 2}

Call a square of the checkerboard “good” if it either has a checker on it orshares a side with a square with a checker on it.

 Say there are $k$ checkers on the board, sitting on squares $s_1, \ldots, s_k$. We can assume WLOG that $s_1, \ldots, s_k$ is ordered in such a way that if $1 \leq i \leq k$, then $s_i$ shares a side with at least one of $s_1, \ldots, s_{i - 1}$. (If such an ordering were impossible, then there would be an $s_i$ which isn’t connected to $s_1$ by a sequence of squares with checkers on them, violating condition (b).)

 Now imagine removing all the checkers and then putting them back onto the squares $s_1, \ldots, s_k$ in that order, counting the number of good squares after eachstep. When we put a checker on $s_1$, we increase the number of good squaresby at most 5: the square $s_1$ becomes good, and each square itshares a side with (at most 4) becomes good. When we put a checker onto $s_i$,where $i > 1$, we increase the number of good squares by at most $3$, since $s_i$ isalready good, and at least $1$ of its at most $4$ neighbors already has a checker on it. 

 So the total number of good squares when we’ve put back every checker is at most $5 + 3(k - 1) = 3k + 2$. In order to satisfy condition (a), we have to make each of the $n^2$ squares good. In other words, we need $3k + 2 \geq n^2$, or 
 \[k \geq \frac{n^2 - 2}{3}.\]

 


	\item (\textit{USAMO 1999}) Let $ABCD$ be a cyclic quadrilateral. Prove that 
 \[|AB - CD| + |AD - BC| \geq 2|AC - BD|.\]


 



\textbf{Solution}

Let arc $AB$ of the circumscribed circle (which we assume WLOG has radius 0.5) have value $2x$, $BC$ have $2y$, $CD$ have $2z$, and $DA$ have $2w$. Then our inequality reduces to, for $x+y+z+w = 180^\circ$: 
 \[|\sin x - \sin z| + |\sin y - \sin w| \ge 2|\sin (x+y) - \sin (y+z)|.\]

 This is equivalent to by sum-to-product and use of $\cos x = \sin (90^\circ - x)$:

 
 \[|\sin \frac{x-z}{2} \sin \frac{y+w}{2}| + |\sin \frac{y-w}{2} \sin \frac{x+z}{2}| \ge 2|\sin \frac{x-z}{2} \sin \frac{y-w}{2}|.\]

 Clearly $90^\circ \ge \frac{x+z}{2} > \frac{x-z}{2} \ge 0^\circ$. As sine is increasing over $[0, \pi/2]$, $|\sin \frac{x+z}{2}| > |\sin \frac{x-z}{2}|$.

 Similarly, $|\sin \frac{y+w}{2}| > |\sin \frac{y-w}{2}|$. The result now follows after multiplying the first inequality by $|\sin \frac{x-z}{2}|$, the second by $|\sin \frac{y-w}{2}|$, and adding. (Equality holds if and only if $x=z$ and $y=w$, ie. $ABCD$ is a parallelogram.)

 --\href{https://artofproblemsolving.com/wiki/index.php?title=User:Suli&action=edit&redlink=1}{Suli} 11:23, 5 October 2014 (EDT)

 


	\item (\textit{USAMO 1999}) Let $p > 2$ be a prime and let $a,b,c,d$ be integers not divisible by $p$, such that

 \[\left\{ \dfrac{ra}{p} \right\} + \left\{ \dfrac{rb}{p} \right\} + \left\{ \dfrac{rc}{p} \right\} + \left\{ \dfrac{rd}{p} \right\} = 2\]
for any integer $r$ not divisible by $p$. Prove that at least two of the numbers $a+b$, $a+c$, $a+d$, $b+c$, $b+d$, $c+d$ are divisible by $p$.
(Note: $\{x\} = x - \lfloor x \rfloor$ denotes the fractional part of $x$.)


 



\textbf{Solution}

We see that $\biggl\{\frac{ra+rb+rc+rd}{p}\biggr\}=0$ means that $p|r(a+b+c+d)$. Now, since $p$ does not divide $r$ and $p$ is prime, their GCD is 1 so $p\mathrel{|}a+b+c+d$.

 Since $\biggl\{\frac{ra}{p}\biggr\}+\biggl\{\frac{rb}{p}\biggr\}+\biggl\{\frac{rc}{p}\biggr\}+\biggl\{\frac{rd}{p}\biggr\}=2$, then we see that they have to represent mods $\bmod\medspace p$, and thus, our possible values of $p$ are all such that $k^4 \equiv 1\pmod{p}$  for all $k$ that are relatively prime to $p$. This happens when $p=3$ or $5$. 

 When $p=3$ then $r$ is not divisible by 3, thus two are $1$, and the other two are $2$. Thus, four pairwise sums sum to 3.

 When $p=5$ then $r$ is not divisible by 5 so $a, b, c, d$ are $1, 2, 3,$ and $4$, so two pairwise sums sum to 5.

 All three possible cases work so we are done.

 (This solution makes absolutely no sense. Why is $k^4\equiv 1$? And how do we know that only $3$ and $5$ work!?)why can't you have a,b,c,d = 2,2,2,4 p = 5 r=1? it doesn't say distinct integers

 


	\item (\textit{USAMO 1999}) Let $a_{1}, a_{2}, \dots, a_{n}$ ($n > 3$) be real numbers such that 
 \[a_{1} + a_{2} + \cdots + a_{n} \geq n \qquad \mbox{and} \qquad a_{1}^{2} + a_{2}^{2} + \cdots + a_{n}^{2} \geq n^{2}.\] Prove that $\max(a_{1}, a_{2}, \dots, a_{n}) \geq 2$.


 



\textbf{Solution 1}

First, suppose all the $a_i$ are positive.  Then
 \[\max(a_1, \dotsc, a_n) \ge \sqrt{\frac{a_1^2 + \dotsb + a_n^2}{n}} \ge \sqrt{n} \ge 2 .\]Suppose, on the other hand, that without loss of generality,
 \[a_1 \ge a_2 \ge \dotsb \ge a_k \ge 0 > a_{k+1} \ge \dotsb \ge a_n,\]with $1\le k <n$.  If $a_1 >2$ we are done, so suppose that $a_1 \le2$.  Then$\sum_{i=1}^k a_i \le 2k$, so 
 \[\sum_{i=k+1}^n -a_i \le 2k-n .\]Since $-a_i$ is a positive real for all $k+1 \le i \le n$, it follows that
 \[\sum_{i=k+1}^n a_i^2 \le \left( \sum_{i=k+1}^n-a_i \right)^2 \le (2k-n)^2 .\]Then
 \begin{align*} \max(a_1, \dotsc, a_n)^2 &\ge \sum_{i=1}^k a_i^2 /k \\ &\ge \left( n^2 - \sum_{i=k+1}^n a_i^2 \right)/k \\ &\ge \frac{n^2 - (2k-n)^2}{k} = 4(n-k). \end{align*}
Since $k<n$, $4(n-k) > 4$.  It follows that $\max(a_1, \dotsc, a_n) \ge \sqrt{4} = 2$, as desired.  $\blacksquare$

 



\textbf{Solution 2}

Assume the contrary and suppose each $a_i$ is less than 2.

 Without loss of generality let $a_1 \le a_2 \le a_3 \le ... \le a_n < 2$, and let $k$ be the largest integer such that $a_k \le 0$ and $0 \le a_{k+1}$ if it exists, or 0 if all the $a_i$ are non-negative. If $k = 0$, then (as $n \ge 4$) $\sum_{i=1}^n a_i^2 < \sum_{i=1}^n 4 = 4n \le n^2$, a contradiction. Hence, assume $n \ge k \ge 1$. Then
 \[\sum_{i=1}^k a_i \ge n - \sum_{i=k+1}^n a_i > n - 2(n-k) = 2k - n.\]Because $a_i \le 0$ for $i \le k$, both sides of the inequality are non-positive, so squaring flips the sign. But we also know that $a_i a_j \ge 0$ for $i, j \le k$, so
 \[\sum_{i=1}^k a_i^2 \le \left(\sum_{i=1}^k a_i \right)^2 < (2k - n)^2 = 4k^2 - 4kn + n^2,\]which results in
 \[\sum_{i=1}^n a_i^2 = \sum_{i=1}^k a_i^2 + \sum_{i=k+1}^n a_i^2 < 4k^2 - 4kn + n^2 + 4(n-k) = n^2 - 4(k-1)(n-k) \le n^2,\]a contradiction to our given condition. The proof is complete.

 



\textbf{Solution 3}

Assume the contrary for maximum $M$ that $M < 2$. Then, all $a_i \leq M < 2$ hence $a_i < 2$. Note that $(a_i)(a_i - 2) < 0$ and therefore $a_i^2 < 2a_i$. When summing over all $n$, that implies $n^2 < 2\sum a_i$. Note that $\sum a_i < 2n$ and therefore $n^2 < 4n$ and thus $n < 4$. However $n \in \mathbb{Z}^+ \neq 1,2,3$ hence a contradiction arises, and so the maximum $M \ge 2$ as requested. $\square$

 
<i>Alternate solutions are always welcome. If you have a different, elegant solution to this problem, please \href{https://artofproblemsolving.com/wiki/index.php?title=1999\_USAMO\_Problems/Problem\_4&action=edit}{add it} to this page.</i>

 


	\item (\textit{USAMO 1999}) The Y2K Game is played on a $1 \times 2000$ grid as follows. Two players in turn write either an S or an O in an empty square. The first player who produces three consecutive boxes that spell SOS wins. If all boxes are filled without producing SOS then the game is a draw. Prove that the second player has a winning strategy.


 



\textbf{Solution}

Label the squares $1$ through $2000$. The first player writes a letter in any square $k$. The second player now selects any empty square $l$ such that $\text{min}\,(|l-1|,|l-2000|,|l-k|) > 3$, and fills in a S. The first player may again play anywhere, say $m$. If the second player can now complete SOS, he does so; otherwise, if $l > m$ he fills an S in $l + 3$ and if $l < m$ he fills an S in $l-3$. 

 We thus have two Ss separated by two empty squares. Note that if any player plays any letter between the two Ss, the other player can always complete a SOS. 

 Now the second player repeats the following algorithm:

 
                \begin{itemize}
                    \item If he can complete SOS on his turn, he does so.\item Otherwise, since there are an odd number of empty squares at the beginning of his turn, there must be either one empty square surrounded by two filled squares or one empty square surrounded by two empty squares. In both cases, the player places an O in that square. In both situations, it is easy to verify that the first player cannot win on the next turn.
                \end{itemize}


                Eventually, assuming the second player has not won yet (and since the first player does not have a winning opportunity yet), this algorithm must stop. It follows that all of the empty squares are now in pairs. Since there are an even number of empty squares, it is the first player's turn. If the first player plays on one of any pair of empty squares that does not have Ss on both its sides, then the second player plays any letter in the other square of the pair. Then, we will only have pairs of empty squares with Ss on both sides remaining (since from the first moves we have constructed at least one, and no one has played between them otherwise the game will have ended), and the first player will have to play between them, so that the second player will win.  

 


	\item (\textit{USAMO 1999}) Let $ABCD$ be an isosceles trapezoid with $AB \parallel CD$. The inscribed circle $\omega$ of triangle $BCD$ meets $CD$ at $E$. Let $F$ be a point on the (internal) angle bisector of $\angle DAC$ such that $EF \perp CD$. Let the circumscribed circle of triangle $ACF$ meet line $CD$ at $C$ and $G$. Prove that the triangle $AFG$ is isosceles.


 



\textbf{Solution}

Quadrilateral $ABCD$ is cyclic since it is an isosceles trapezoid. $AD=BC$. Triangle $ADC$ and triangle $BCD$ are reflections of each other with respect to diameter which is perpendicular to $AB$. Let the incircle of triangle $ADC$ touch $DC$ at $K$. The reflection implies that $DK=CE$, which then implies that the excircle of triangle $ADC$ is tangent to $DC$ at $E$. Since $EF$ is perpendicular to $DC$ which is tangent to the excircle, this implies that $EF$ passes through center of excircle of triangle $ADC$.

 We know that the center of the excircle lies on the angular bisector of $DAC$ and the perpendicular line from $DC$ to $E$. This implies that $F$ is the center of the excircle. 

 Now $\angle GFA=\angle GCA=\angle DCA$. $\angle ACF=90^\circ+\frac{\angle DCA}{2}$.This means that $\angle AGF=90^\circ-\frac{\angle ACD}{2}$. (due to cyclic quadilateral $ACFG$ as given).Now $\angle FAG =180^\circ- (\angle AFG + \angle FGA)=90^\circ-\frac{\angle ACD}{2}=\angle AGF$.

 Therefore $\angle FAG=\angle AGF$. QED.

 


\end{enumerate}